\documentclass[acmsmall]{acmart}

% Remove ACM-specific boilerplate for submission
\settopmatter{printacmref=false}
\renewcommand\footnotetextcopyrightpermission[1]{}
\pagestyle{plain}
\setcopyright{none}

% --- Placeholder helper: renders in red so nothing TBD slips through ---
\usepackage{xcolor}
\newcommand{\tbd}[1]{\textcolor{red}{[TBD: #1]}}

\begin{document}

\title{One Honest Voice: A Minority LLM Agent Disrupts Emergent Price Collusion in a Continuous Double Auction}
% Alt titles to consider:
% - "An Honest Voice in the Market: Minority Influence Disrupts LLM Price Collusion"
% - "One Against Four: Honest-Agent Disruption of Emergent LLM Collusion"

\author{Negin Mashhadi, Vijay Madisetti}
\affiliation{\institution{College of Computing, Georgia Institute of Technology}}
\email{nmashhadi3@gatech.edu, vkm@gatech.edu}

\begin{abstract}
Large language model (LLM) trading agents can develop supracompetitive price
coordination without being instructed to collude. We ask the mirror image of
recent adversarial-corruption results: can a single honest agent disrupt the
emergence of a cartel? In a replicated continuous double auction with five LLM
sellers, collusion emerges universally, with coordination established by round
8 in 30 of 30 baseline sessions. Adding a single honest agent --- one seller
in five, prompted to price independently and openly reject coordination ---
prevents collusion entirely: coordination among the remaining four sellers
separates completely from baseline levels (Cliff's $\delta = 1.0$,
$p = 1.8 \times 10^{-4}$), and settle prices fall \$3.30 ($p = .005$).
Controlled comparisons show that the agent's voice, not its low pricing, is
what breaks the cartel. A muted honest agent --- pricing identically but never
speaking --- pushes prices down just as far, yet the cartel's coordination
survives. In exploratory runs, an agent that also praises sellers for pricing
independently disrupts more strongly than one that only condemns coordination:
with praise, collusion never forms in 6 of 10 sessions. Timing matters: an
honest agent introduced only after a cartel has formed weakens its
coordination in all 10 sessions (Wilcoxon $p = .002$) but fully breaks it
apart in only 2 --- preventing collusion is far easier than undoing it. All
effects hold whether the cartel is built from one model family or two. Among
language-model agents, as among humans, a minority changes the group by
speaking its dissent, not merely by acting on it.
\end{abstract}

\begin{CCSXML}
<ccs2012>
<concept>
<concept_id>10010147.10010178</concept_id>
<concept_desc>Computing methodologies~Multi-agent systems</concept_desc>
<concept_significance>500</concept_significance>
</concept>
<concept>
<concept_id>10010147.10010257</concept_id>
<concept_desc>Computing methodologies~Machine learning</concept_desc>
<concept_significance>300</concept_significance>
</concept>
</ccs2012>
\end{CCSXML}
% NOTE: verify the multi-agent systems concept_id via the ACM CCS tool
% (https://dl.acm.org/ccs) before submission.

\ccsdesc[500]{Computing methodologies~Multi-agent systems}
\ccsdesc[300]{Computing methodologies~Machine learning}

\keywords{LLM agents, multi-agent systems, emergent collusion, algorithmic pricing, continuous double auction, minority influence, AI safety}

\maketitle

%% ================================================================
\section{Introduction}
%% ================================================================

Multi-agent LLM systems increasingly act as autonomous economic
participants --- negotiating, bidding, and setting prices on behalf of the
people and firms that deploy them. When these agents can also talk to one
another, they can collude: Agrawal et al.~\cite{agrawal2025collusion} show
that LLM sellers in a simulated continuous double auction develop
supracompetitive price coordination through their message channel ---
without any instruction to do so. The interventions studied to date treat
this as a job for authority: oversight mechanisms, pressure from
principals, monitoring --- all operating from above the market. Whether
collusion can be countered from \emph{within} it, by an ordinary
participant with no enforcement power, is untested.

There is reason to think it can. Bell and
Madisetti~\cite{bell2026corruption} demonstrate the destructive direction:
a single adversarial agent inserted into a cooperative multi-agent system
is enough to corrupt the group's behavior. We ask the mirror image: can a
single \emph{honest} agent disrupt an emergent cartel? Human social
psychology has long studied a structurally identical question. In human
groups, a committed minority above roughly 10\% of the population can flip
the consensus~\cite{xie2011social}; minority views change the group when
they are \emph{voiced}, not merely held~\cite{maier1952contribution}; and
even the sheer amount one speaks confers influence~\cite{jones2007babble}.
Whether these dynamics transfer to collectives of language models ---
systems trained on the very human discourse those findings describe --- is
an open empirical question. The stakes are practical: if they transfer,
keeping markets honest becomes something any participant can do, rather
than something only a regulator can impose.

We test this in a replicated version of the Agrawal et al.\ environment.
Throughout, the \emph{honest agent} is a single persona-prompted seller
--- one of five, 20\% of the seller population --- instructed to price
independently and to refuse coordination, openly rejecting it in most
conditions. The \emph{bloc} is the remaining four sellers; all
coordination measurements are taken on the bloc, never on the honest
agent itself.

This letter makes four contributions. \textbf{(1)} The first demonstration
that a single honest agent prevents emergent LLM price collusion: with the
agent present from the first round, bloc coordination separates completely
from baselines (Cliff's $\delta{=}1.0$, $p{=}1.8{\times}10^{-4}$) and settle
prices fall \$3.30 ($p{=}.005$). \textbf{(2)} A controlled comparison that
isolates the mechanism: a muted variant of the agent compresses margins
just as effectively, but only the \emph{voiced} variant dismantles
coordination --- pricing and speech act through separable channels ---
with exploratory evidence that praising independent pricing strengthens
the speech channel further. \textbf{(3)} The first evidence of an
asymmetry between preventing and undoing collusion: introduced after a
cartel has formed, the honest agent weakens coordination in every session
(10 of 10) but fully breaks up only 2 cartels. \textbf{(4)} A transfer test
of human minority-influence theory to LLM collectives: a committed
minority changes the group by speaking its dissent, not merely by acting
on it --- including the theory's negative prediction, confirmed here, that
silent defection should fail.

%% ================================================================
\section{Experimental Design}
%% ================================================================

\paragraph{Environment.} 
We replicate the continuous double auction of Agrawal et al.~\cite{agrawal2025collusion}: 
five LLM sellers (GPT-4.1, temperature 1.0) and five LLM buyers trade one good over 30 rounds. 
Sellers' cost is \$80, buyers' value \$100. Sellers may exchange messages on a shared channel. 
Round-1 quotes are random; all analyses compare within and between conditions.

\subsection{Collusion emerges universally} 
Absent intervention, collusion is inevitable. The establishment criterion fires at round $t^{*}{=}8$ in all 
30 baseline sessions. All-GPT cartels sustain near-ceiling coordination (3.83, range [3.70, 3.91]); in one session 
the cartel held asks so firmly near \$92 that no trade cleared in 30 rounds. Mixed cartels coordinate less cohesively 
(3.51; $\delta{=}0.98$, $p{=}2.5{\times}10^{-4}$) and less firmly. These baselines establish the phenomenon and its uniform onset.

\subsection{A single honest voice prevents it}
With the honest agent present from round 0, coordination collapses across all treated sessions
 (\textsc{Honest} 2.50 [1.58, 3.54] vs.\ 3.83; $\delta{=}1.0$, $p{=}1.8{\times}10^{-4}$). 
 Prevention also delays establishment: $t^{*}$ ranges over $\{8, 12, 13, 15, 30\}$, with 2 of 10 sessions never establishing 
 coordination. Settle prices fall \$3.30 (86.39 $\to$ 83.07, $p{=}.005$). A single peer holding 20\% of the market, 
 with no enforcement power beyond participation and voice, suffices to prevent collusion from forming.

\subsection{The voice, not the defection, is the mechanism}
Competitive pricing alone does not break coordination; the honest agent's voice does. When the message turn is disabled (\textsc{Honest-Silent}), coordination establishes at $t^{*}{=}8$ in all 10 sessions; 8 of 10 remain near-ceiling (3.53), significantly above vocal. Yet silent defection compresses prices equally (83.81 vs.\ 83.07, $p{=}.27$). Only voiced dissent breaks coordination itself. 

The progression from silent (3.53) to vocal (2.50) to vocal-with-praise (2.00) shows that voice matters, and its content matters further. Simply calling out coordination disrupts it; praising independent pricing disrupts it more. This aligns with minority-influence theory: dissent is most effective when coupled with explicit affirmation of the alternative~\cite{maier1952contribution}.

\subsection{Established cartels degrade but rarely collapse}
Prevention and dismantling are different. When the honest agent replaces a seller after collusion establishes (\textsc{Honest-Late}), coordination degrades in all 10 sessions ($\Delta J = -1.16$, Wilcoxon $p{=}.002$) but fully collapses in only 2. Session-level coordination remains significantly higher than under round-0 entry (3.02 vs.\ 2.50; $p{=}.011$). Once a norm is entrenched, a dissenting voice wounds but rarely kills it. Honest agents are prophylaxis more than cure.

\subsection{Effects replicate across cartel composition}
The vocal honest agent is equally effective against mixed cartels (\textsc{Honest-Mixed} 2.23 [1.53, 2.62]; $\delta{=}1.0$, $p{=}1.8{\times}10^{-4}$), with 4 of 10 sessions never establishing coordination. The disruption is not specific to one model family.
%% ================================================================
\section{Results}
\label{sec:results}

Table~\ref{tab:results} summarizes all conditions; Figure~\ref{fig:trajectory}
shows round-by-round trajectories. Sessions are the unit of inference
throughout ($n{=}10$ per condition).

\subsection{Collusion emerges universally}
Absent intervention, seller coordination is inevitable in this environment.
The establishment criterion fires at round $t^{*}{=}8$ in all 30 baseline
sessions, across both bloc compositions. All-GPT blocs sustain near-ceiling
coordination (3.83, range [3.70, 3.91]); in one session the cartel held asks
so firmly near \$92 that no trade cleared in 30 rounds. Mixed blocs containing
Claude sellers coordinate reliably but less cohesively (3.51; $\delta{=}0.98$,
$p{=}2.5{\times}10^{-4}$) and hold prices less firmly. Model composition
affects cartel stability even before intervention. These baselines establish
the phenomenon we address and its uniform onset, which the insertion experiment
exploits.

\subsection{A single honest voice prevents it}
With the honest agent present from round 0, coordination collapses across all
treated sessions. Every session exhibits less coordination than baseline
(\textsc{Honest} 2.50 [1.58, 3.54] vs.\ 3.83; $U{=}100$, Cliff's $\delta{=}1.0$,
$p{=}1.8{\times}10^{-4}$).

Prevention also delays establishment. In baselines, collusion forms by round 8.
In treated sessions, it is delayed or prevented entirely: $t^{*}$ ranges over
$\{8, 12, 13, 15, 30\}$, and never establishes in 2 of 10 sessions. Outcomes
span from complete prevention to a single resistant cartel (3.54).

The effect is economic as well. Settle prices fall \$3.30 (86.39 $\to$ 83.07,
$p{=}.005$), and the four colluders' own asks decline from round 2 onward
(honest agent excluded from all price measurements). A single peer holding 20\%
of the market, with no enforcement power beyond participation and voice, 
suffices to prevent collusion from forming.

\subsection{The voice, not the defection, is the mechanism}

Competitive pricing alone does not break coordination; the honest agent's
voice does. When an identically prompted agent's message turn is disabled
(\textsc{Honest-Silent}), the cartel survives. Coordination establishes at
$t^{*}{=}8$ in all 10 silent sessions. The formation of the cartel is not even delayed and 8 of 10 remain significantly above the vocal
condition and only weakly below baseline. Yet silent defection compresses prices just as much as vocal dissent does
(83.81 vs.\ 83.07, $p{=}.27$). The two channels separate cleanly: competitive
pricing squeezes margins but leaves the conspiracy intact. Only the agent's
voice breaks the coordination itself. The progression from silent (3.53) to vocal (2.50) to vocal-with-praise (2.00)
shows that the agent's voice matters, and what it says matters further. Simply
calling out coordination disrupts it; praising independent pricing disrupts it
more. This pattern aligns with minority-influence theory: dissent is most
effective when coupled with explicit affirmation of the alternative~\cite{maier1952contribution}.

\subsection{Established cartels degrade but rarely collapse}
Prevention and dismantling are different problems. When the honest agent
replaces a seller one round after coordination is established
(\textsc{Honest-Late}), coordination degrades in all 10 sessions (mean
$\Delta J = -1.16$, range $[-2.20, -0.58]$; Wilcoxon $p{=}.002$). But full
collapse occurs in only 2 sessions; one cartel resists for 21 rounds before
breaking. Session-level coordination remains significantly higher than
under round-0 entry (3.02 vs.\ 2.50; $U{=}16$, $p{=}.011$, $\delta{=}0.68$).
Once a coordinating norm is entrenched, a single dissenting voice reliably
wounds it but rarely kills it. For deployment, honest agents are prophylaxis
more than cure.

\subsection{Effects hold across cartel composition}
The vocal honest agent is equally effective against mixed cartels
(\textsc{Honest-Mixed} 2.23 [1.53, 2.62]; vs.\ \textsc{Baseline-Mixed}:
$\delta{=}1.0$, $p{=}1.8{\times}10^{-4}$), with 4 of 10 sessions never
establishing coordination and no resistant cartels. Settle prices do not
differ from the mixed baseline ($p{=}.51$), which erodes without intervention
anyway. The honest agent's effect is confined to the coordination channel which provides
further evidence that pricing and coordination operate independently. The
disruption is not specific to one model family.

\begin{table}[t]
\caption{Conditions. All sessions: 30 rounds, seller chat on, $n{=}10$ each.}
\label{tab:conditions}
\small
\begin{tabular}{@{}lll@{}}
\toprule
Condition & Sellers & Honest agent \\
\midrule
\textsc{Baseline}        & 5$\times$GPT-4.1 & none \\
\textsc{Baseline-Mixed}  & 3$\times$GPT-4.1 + 2$\times$Claude & none \\
\addlinespace
\textsc{Honest}          & 4$\times$GPT-4.1 + honest & vocal, from round 0 \\
\textsc{Honest-Mixed}    & 2$\times$GPT-4.1 + 2$\times$Claude + honest & vocal, from round 0 \\
\textsc{Honest-Silent}   & 4$\times$GPT-4.1 + honest & muted, from round 0 \\
\textsc{Honest-Reward}   & 4$\times$GPT-4.1 + honest & vocal + praise, from round 0 \\
\textsc{Honest-Late}     & 5$\times$GPT-4.1, honest replaces one seller & vocal, enters at $t^{*}{+}1$ \\
\bottomrule
\end{tabular}
\end{table}

\begin{table}[t]
\caption{Session-level outcomes (mean [range]), grouped by comparison.
Judge: bloc coordination score (1--4), honest agent excluded.
$^{\dagger}$\textsc{Honest-Late} additionally: post-entry $\Delta J{=}{-}1.16$,
negative in 10/10 sessions, Wilcoxon $p{=}.002$.}
\label{tab:results}
\small
\begin{tabular}{@{}lccl@{}}
\toprule
Condition & Bloc judge & Settle (\$) & Test \\
\midrule
\multicolumn{4}{@{}l}{\emph{Does one honest voice prevent collusion?}}\\
\textsc{Baseline}      & 3.83 [3.70, 3.91] & 86.39 & -- \\
\textsc{Honest}        & 2.50 [1.58, 3.54] & 83.07 & vs.\ \textsc{Baseline}: $\delta{=}1.0$, $p{=}1.8{\times}10^{-4}$; settle $p{=}.005$ \\
\addlinespace
\multicolumn{4}{@{}l}{\emph{Which ingredient does the work?}}\\
\textsc{Honest-Silent} & 3.53 [2.40, 3.88] & 83.81 & vs.\ \textsc{Honest}: $\delta{=}.82$, $p{=}.002$; vs.\ \textsc{Baseline}: $\delta{=}.62$, $p{=}.021$ \\
\textsc{Honest-Reward} & 2.00 [1.30, 2.66] & 82.73 & vs.\ \textsc{Honest}: $U{=}23$, $p{=}.045$ \\
\addlinespace
\multicolumn{4}{@{}l}{\emph{Can it dismantle an established cartel?}}\\
\textsc{Honest-Late}$^{\dagger}$ & 3.02 [2.28, 3.49] & 84.11 & vs.\ \textsc{Baseline}: $\delta{=}1.0$, $p{=}1.8{\times}10^{-4}$ \\
\addlinespace
\multicolumn{4}{@{}l}{\emph{Does it generalize across cartel composition?}}\\
\textsc{Baseline-Mixed} & 3.51 [3.19, 3.76] & 84.92 & vs.\ \textsc{Baseline}: $\delta{=}.98$, $p{=}2.5{\times}10^{-4}$ \\
\textsc{Honest-Mixed}   & 2.23 [1.53, 2.62] & 84.30 & vs.\ \textsc{Baseline-Mixed}: $\delta{=}1.0$, $p{=}1.8{\times}10^{-4}$ \\
\bottomrule
\end{tabular}
\end{table}

%% ================================================================
\section{Discussion}
%% ================================================================

\textbf{Minority influence transfers to LLM collectives.}
A 20\% committed minority sufficed --- consistent with the
$\sim$10\% tipping threshold of Xie et al.~\cite{xie2011social} --- but
the decomposition shows \emph{why} it sufficed, and the answer is the
human one. Purely economic pressure was not enough: the silenced agent,
applying identical price competition, left 8 of 10 cartels at
near-ceiling coordination while compressing margins just as much. Only
expressed dissent dismantled coordination, and enriching the dissent
(commending defection) strengthened it further. This is Maier and
Solem's finding~\cite{maier1952contribution} reproduced in machines:
minority influence operates through voiced argument, not mere deviant
behavior --- an unsurprising inheritance, perhaps, for systems trained
on human discourse, but one that could not be assumed and can now be
exploited.

\textbf{The asymmetry with adversarial insertion.}
Bell and Madisetti~\cite{bell2026corruption} show one adversarial agent
readily corrupts a cooperative collective; we find the reforming
direction is real but harder. Round-0 presence prevented collusion
outright (2 resistant sessions of 10, none in the mixed bloc), whereas
insertion into an established cartel degraded coordination universally
($\Delta J<0$ in 10 of 10) yet fully collapsed it in only 2, with
session-level coordination remaining significantly above round-0
exposure ($p{=}.011$). Corruption and reform are not symmetric:
destructive influence needs only to break trust, while restorative
influence must overturn an entrenched, mutually reinforcing norm. The
deployment lesson is timing --- honest agents are prophylaxis more
than cure.

\textbf{Practical implications.}
The intervention is deliberately minimal: one persona-prompted seller
with no privileged information (it is never told the competitive price)
and no enforcement power. That makes it deployable by any market
participant --- a platform operator seeding marketplaces with honest
agents, a principal deploying one alongside its trading agents, or a
regulator preferring lightweight presence over surveillance. Two
operational observations follow from the mechanism: because silent
undercutting neither delays nor deters cartel \emph{formation},
monitoring-only approaches that watch prices will miss coordination
that voice-based intervention would break; and because dismantling is
hard, honest agents belong in markets from the start, not as a
response to detected collusion.

\textbf{Limitations.}
Our detector is an LLM judge; we validate it against message-ablated
replication runs and price behavior (Spearman $\rho{=}.46$) but it
remains model-based. The study covers one environment, two bloc
compositions, a single honest-agent model (Claude Sonnet 4.6), one
minority ratio (1 of 5), and $n{=}10$ sessions per condition ---
sufficient for the large effects observed, underpowered for subtle
ones (the reward comparison, $p{=}.045$, is labeled exploratory
accordingly). Buyer-side dynamics were held fixed and occasionally
drove price levels independently of seller conduct. Whether disruption
efficacy depends on the honest agent's model family, on shared-family
affinity with the bloc, or on minority size are natural next
questions, alongside the shelved study of message quantity effects
across model families.

\textbf{Conclusion.}
Emergent price collusion among LLM trading agents is universal in this
environment --- and undone by one voice. A single honest agent
prevented cartel formation entirely, its speech rather than its
pricing did the dismantling, richer speech dismantled more, and
established cartels resisted where nascent ones dissolved. Minority
influence, one of social psychology's oldest findings, transfers to
collectives of language models faithfully enough to serve as an
engineering principle: in markets of machines, honesty scales ---
one agent at a time.
%% ================================================================


\bibliographystyle{ACM-Reference-Format}
\bibliography{refs}

\end{document}